% SEGMENT OLP-0526-S01
% Part: counterfactuals
% Chapter: minimal-change-semantics
% Section: antecedent-strengthening

\documentclass[../../../include/open-logic-section]{subfiles}

\begin{document}

\olfileid{cnt}{min}{agg}

\olsection{முன்னிலையை வலுப்படுத்துதல்}

``முன்னிலையை வலுப்படுத்துதல்'' என்பது $!A \lif !C \Entails (!A
\land !B) \lif !C$ என்ற அனுமானத்தைக் குறிக்கிறது. இது பொருள்சார்
நிபந்தனைக்குச் செல்லுபடியாகும்; ஆனால் உண்மைக்கு மாறான நிபந்தனைகளுக்குச்
செல்லுபடியாகாது. இந்தத் தீக்குச்சியை நான் உரசியிருந்தால் அது பற்றியிருக்கும்
என்பது மெய் எனக் கொள்க. (அதாவது, தீக்குச்சியிலோ தீப்பெட்டியின் உரசும்
பரப்பிலோ குறை ஏதுமில்லை; நான் தீக்குச்சியை முறிக்கமாட்டேன்; இப்படியே
பல.) ஆனால் விண்வெளியில் இந்தத் தீக்குச்சியை நான் உரசியிருந்தால் அது
பற்றியிருக்கும் என்பது மெய்யல்ல. எனவே பின்வரும் அனுமானம் செல்லுபடியாகாது:
% SOURCE CORRECTION TA-CF-006: the strengthened antecedent is striking, not lighting, the match in outer space.
\begin{quote}
  தீக்குச்சி உரசப்பட்டிருந்தால், அது பற்றியிருக்கும்.

  எனவே, தீக்குச்சி விண்வெளியில் உரசப்பட்டிருந்தால், அது பற்றியிருக்கும்.
\end{quote}

நிபந்தனைக் கூற்றுகள் பற்றிய லூயிஸ்--ஸ்டால்நேக்கர் விளக்கம் இதை
விளக்குகிறது: தீக்குச்சியை நான் விண்வெளியில் உரசும் மிக நெருங்கிய உலகம்,
தீக்குச்சியை நான் உரசும் மிக நெருங்கிய உலகைவிட நடப்பு உலகிலிருந்து மிக
அதிகத் தொலைவில் உள்ளது. ஆகவே பிந்தைய உலகில் தீக்குச்சி பற்றுவது மெய்யாக
இருந்தாலும், முந்தைய உலகில் அது மெய்யல்ல. அவ்வாறே இருக்க வேண்டும்.

\begin{ex}
  கோளப் பொருண்மையியல் இந்த அனுமானத்தைச் செல்லாததாக்குகிறது; அதாவது $p
  \cif r \Entails/ (p \land q) \cif r$. $W = \{w, w_1, w_2\}$,
  $O_w = \{\{w\}, \{w, w_1\}, \{w, w_1, w_2\}\}$, $V(p) = \{w_1,
  w_2\}$, $V(q) = \{w_2\}$, $V(r) = \{w_1\}$ ஆகியவற்றைக் கொண்ட
  மாதிரி $\mModel{M} = \tuple{W, O, V}$ ஐக் கருதுக. $p$-ஐ ஏற்கும்
  கோளம் $S = \{w, w_1\}$ ஒன்று உள்ளது; அதிலுள்ள எல்லா உலகங்களிலும்
  $p \lif r$ மெய்; எனவே $\mSat{M}{p \cif r}[w]$. மேலும், $(p \land
  q)$-ஐ ஏற்கும் கோளம் $S' = \{w, w_1, w_2\}$ ஒன்று உள்ளது. ஆனால்
  $\mSat/{M}{(p \land q) \lif r}[w_2]$; எனவே $\mSat/{M}{(p \land q)
  \cif r}[w]$ (\olref{fig:antecedent-strengthening} ஐப் பார்க்கவும்).

\begin{figure}
\begin{center}
\begin{tikzpicture}[layerwidth=1.5,scale=.6]\tiny
  \spheresystem{3}
  \propositionintersect{-10}{3}{30}{5}
  \begin{scope}
    \clip plot[smooth,tension=1] coordinates
          {(0,4.5) (30:.95) (4.5,-3.5)} -- (4.5,4.5) ;
    \spherefill{2}
  \end{scope}
  \draw[proposition] plot[dashed,smooth,tension=1] coordinates
          {(0,4.5) (30:.95) (4.5,-3.5)} ;
  \proposition{30}{2}{30}{5}
  \path (0,0) node[world] {$w$};
  \spherepos{30}{2}{node[world] {$w_1$}}
  \spherepos{-10}{3}{node[world] {$w_2$}}
  \spherepos{-10}{4}{node {$q$}}
  \spherepos{30}{4}{node {$r$}}
  \draw (1,4.5) node[below] {$p$};
\end{tikzpicture}
\caption{முன்னிலையை வலுப்படுத்துதலுக்கான எதிரெடுத்துக்காட்டு}
\ollabel{fig:antecedent-strengthening}
\end{center}
\end{figure}
\end{ex}

\end{document}
